\scanpage{17}
$x$ in question, or the $i$ in question) form a \uline{torsor} under the group $\handwritten{Z^+}=$\deleted{$\handwritten{\unsure{U\cap}}\operatorname{Cent}(u)$}\inserted{$\handwritten{Z\cap U}$}, which is a unipotent group\inserted{\handwritten{\deleted{(in fact, $\unsure{\subset|c}U$ introduced in (iv))}}}\handwritten{.}
\margin{\handwritten{cf below definition of $U$. We have $Z\simeq Z_{\unsure{o|0}}.Z^+$ \deleted{\unsure{Commuted|Commutative}} (semidirect product) with $Z_o=Z\cap L=\operatorname{Cent}(f)$.}}
Finally, still assuming $e\handwritten{\ne}0$, the set of $x$ in question is a torsor under \deleted{$\handwritten{\mathfrak{g}}^{e}=$ (centralizer of $e$ \deleted{\unsure{in}} in $\handwritten{\mathfrak{g}}$)$=\operatorname{Ker}\operatorname{ad}(e)=\operatorname{Lie}\operatorname{Cent}(u)$}\inserted{$\handwritten{\operatorname{Lie}(Z^+)=\mathfrak{g}^{e}\cap\operatorname{Fil}^{1}(\mathfrak{g})}$} (\inserted{\handwritten{torsor}}for addition, i.e. is of the form $a+\handwritten{\mathfrak{g}}^e$, for a suitable $a\handwritten{\in\mathfrak{g}}$). \deleted{(and we have $\handwritten{\mathfrak{g}^{e}\subset\sum_{k\geq1}\mathfrak{g}}(k)=\operatorname{Fil}^{1}(\handwritten{\mathfrak{g}})$)}

(iii) Let $x\inserted{\handwritten{\ne0}}$ be a Dynkin element of $\handwritten{\mathfrak{g}}$, thus associated with a Dynkin homomorphism $i:\underline{G}_m\handwritten{\longrightarrow}G$. Then the elements\inserted{\handwritten{$e$}} of $\handwritten{\mathfrak{g}}$ such that $(x,e)$ satisfies\deleted{nt} the conditions of (i) form a dense open subset\inserted{\handwritten{\deleted{$\unsure{\mathfrak{g}(2)^*|\ne0}$}}} of $\handwritten{\mathfrak{g}}(2)$ (namely the open subset\inserted{\handwritten{$\mathfrak{g}(2)^*$}} of points $e$ such that $\operatorname{ad}(e):\handwritten{\mathfrak{g}}(0)\handwritten{\to\mathfrak{g}}(2)$ is surjective), \deleted{which}\inserted{\unsure{This|This}} is a \deleted{\uline{torsor}}\inserted{\handwritten{homogeneous space}} under\inserted{\handwritten{$L=$}}$\operatorname{Cent}(x)=\operatorname{Cent}(i)$, which is a reductive group.

(iv) For every \deleted{pair $(x,e)$ satisfying the conditions of (i),}\inserted{Dynkin element $x$ of $\handwritten{\mathfrak{g}}$, $x=\operatorname{Lie}(i)(1)$,} let us consider the subgroups $P,U,\inserted{L}$ of $G$ defined by the conditions
\[
\operatorname{Lie}P=\handwritten{\sum_{k\geq0}\mathfrak{g}}(k)\quad,\quad
\operatorname{Lie}U=\handwritten{\sum}_{k\handwritten{\geq}\deleted{0}\handwritten{1}}\handwritten{\mathfrak{g}}(k),\quad
\operatorname{Lie}L=\handwritten{\mathfrak{g}}(0).
\]
Then $P$ is a parabolic subgroup of $G$, $U$ is its unipotent radical, and $L=\operatorname{Cent}(x)=\operatorname{Cent}(i)$ \deleted{\unsure{in|a}} is a Levi subgroup of $G$. If $\deleted{\unsure{a}}u\inserted{=\exp(e)}$ is a unipotent element of $G$ such that $(i,u)$ is a Dynkin pair (i.e. induced by an $f:\unsure{\operatorname{SZ}|\operatorname{Sl}}(2,k)\handwritten{\to}G$), $P$ (thus a\deleted{\unsure{a}}lso $U$) is defined in terms of $u$ (i.e; two $x$ associated with the same $u$ give the same $P$, cf. (ii)); the same holds for the subgroup $H=i(\underline{G}_m).U$ of $G$ (NB $i(\underline{G}_m)$ normalizes $U$, and $H$ is a semidirect product), and for $i$ modulo $U$ i.e. for $\underline{G}_m\handwritten{\longrightarrow}H/U$ induced by $i$. \deleted{A}Moreover, the map
\[
g\handwritten{\longmapsto}\operatorname{ad}(g).x
\]
is an\deleted{e} isomorphism of $U$ with $x+\operatorname{Fil}^{1}\handwritten{\mathfrak{g}}\inserted{\handwritten{=x+\operatorname{Lie}U}}$ (where we set $\operatorname{Fil}^{k}\handwritten{\mathfrak{g}}=\handwritten{\sum_{j\geq k}\mathfrak{g}}(j)$), and the map
\[
g\handwritten{\longmapsto}\operatorname{ad}(g).\handwritten{e}
\]
is a surjective map from $U$ onto $e+\operatorname{Fil}^{\handwritten{3}}(\handwritten{\mathfrak{g}})$, or from $P$ onto $\handwritten{\mathfrak{g}}(2)^{\handwritten{*}}+\operatorname{Fil}^{\handwritten{3}}(\handwritten{\mathfrak{g}})$, where $\handwritten{\mathfrak{g}}(2)^{\handwritten{*}}$ is the open subset of $\handwritten{\mathfrak{g}}(2)$ considered in (iii).
% archive sheet number: 16
